% GENERATED FILE. Edit canonical lessons, then run npm run generate:books.
% canonical-source-hash: fnv1a64:1ae38e231174ad50
% canonical-lessons: GU-C533-paso, GU-C533-copat, GU-C533-dav, GU-C533-baji, GU-C533-khelvu

\chapter{A Roll of the Dice}
\label{ch:gu-a-roll-of-the-dice}
\hlchaptermodality{533}

\begin{chapteropening}
\textbf{By the end of this chapter:} \emph{I can talk about dice games and whose turn it is.}

Complete the last lesson of chapter 533: I can talk about dice games and whose turn it is.
\end{chapteropening}

\section[\texorpdfstring{\gu{પાસો} (pāso)}{પાસો (pāso)}]{\gu{પાસો} (pāso) — a die}

\label{lesson:GU-C533-paso}



\begin{quote}
Before the new one: say the Gujarati for convenient, then the Gujarati for a resignation.
\end{quote}

\subsection*{You'll want to know: \gu{પાસો}}
\textbf{\gu{પાસો}} — \emph{pāso} — \textquotedblleft{}a die\textquotedblright{}.

\begin{quote}
\gu{જો} \gu{તારો} \gu{પાસો} \gu{છ} \gu{પર} \gu{પડશે} \gu{તો} \gu{તું} \gu{ફરી} \gu{રમીશ}.
\end{quote}

\begin{quote}
\emph{jo tāro pāso cha par paḍśe to tũ pharī ramīś.}
\end{quote}

\begin{quote}
If your die lands on six, you'll play again.
\end{quote}

A masculine noun.

\begin{cousinweb}
\textbf{\gu{પાસો}} (\emph{pāso}) goes back to Sanskrit \emph{pāśa}, \textquotedblleft{}a die\textquotedblright{}, and CDIAL lists the Gujarati word (CDIAL 8132).
\end{cousinweb}

\begin{grammarlens}[title={What you've built}]
One more word for saying what happens if.
\end{grammarlens}

\subsection*{Your turn}
\begin{itemize}
\raggedright
  \item \emph{Say it:} \emph{pāso}
  \item \emph{Say it:} \emph{pāso}, once more
  \item \emph{Say it:} \emph{jo tāro pāso cha par paḍśe to tũ pharī ramīś.} — If your die lands on six, you'll play again.
  \item \emph{Recall:} say the Gujarati for convenient, then the Gujarati for a resignation, then say \emph{pāso} again
\end{itemize}

\begin{tcolorbox}[breakable,colback=teal!4,colframe=teal!35!black,title={Before you move on}]
What does \gu{પાસો} mean? (\textquotedblleft{}A die\textquotedblright{}.) Say it once more.
\end{tcolorbox}
\section[\texorpdfstring{\gu{ચોપાટ} (copāṭ)}{ચોપાટ (copāṭ)}]{\gu{ચોપાટ} (copāṭ) — chaupat (a dice board game)}

\label{lesson:GU-C533-copat}



\begin{quote}
Before the new one: say the Gujarati for a resignation, then the Gujarati for a die.
\end{quote}

\subsection*{You'll want to know: \gu{ચોપાટ}}
\textbf{\gu{ચોપાટ}} — \emph{copāṭ} — \textquotedblleft{}chaupat (a dice board game)\textquotedblright{}.

\begin{quote}
\gu{જો} \gu{બહાર} \gu{વરસાદ} \gu{પડશે} \gu{તો} \gu{આપણે} \gu{ઘરમાં} \gu{ચોપાટ} \gu{રમીશું}.
\end{quote}

\begin{quote}
\emph{jo bahār varsād paḍśe to āpaṇe gharmā̃ copāṭ ramīśũ.}
\end{quote}

\begin{quote}
If it rains outside, we'll play chaupat indoors.
\end{quote}

A feminine noun.

\begin{cousinweb}
\textbf{\gu{ચોપાટ}} (\emph{copāṭ}) comes from \emph{*catuṣpaṭṭa}, \textquotedblleft{}a four-sided tablet\textquotedblright{}, a form CDIAL reconstructs, and CDIAL lists the Gujarati word as a dice board (CDIAL 4639).
\end{cousinweb}

\begin{grammarlens}[title={What you've built}]
One more word for saying what happens if.
\end{grammarlens}

\subsection*{Your turn}
\begin{itemize}
\raggedright
  \item \emph{Say it:} \emph{copāṭ}
  \item \emph{Say it:} \emph{copāṭ}, once more
  \item \emph{Say it:} \emph{jo bahār varsād paḍśe to āpaṇe gharmā̃ copāṭ ramīśũ.} — If it rains outside, we'll play chaupat indoors.
  \item \emph{Recall:} say the Gujarati for a resignation, then the Gujarati for a die, then say \emph{copāṭ} again
\end{itemize}

\begin{tcolorbox}[breakable,colback=teal!4,colframe=teal!35!black,title={Before you move on}]
What does \gu{ચોપાટ} mean? (\textquotedblleft{}Chaupat (a dice board game)\textquotedblright{}.) Say it once more.
\end{tcolorbox}
\section[\texorpdfstring{\gu{દાવ} (dāv)}{દાવ (dāv)}]{\gu{દાવ} (dāv) — a turn (in a game)}

\label{lesson:GU-C533-dav}



\begin{quote}
Before the new one: say the Gujarati for a die, then the Gujarati for chaupat.
\end{quote}

\subsection*{You'll want to know: \gu{દાવ}}
\textbf{\gu{દાવ}} — \emph{dāv} — \textquotedblleft{}a turn (in a game)\textquotedblright{}.

\begin{quote}
\gu{જો} \gu{હવે} \gu{મારો} \gu{દાવ} \gu{હોત} \gu{તો} \gu{હું} \gu{જીતી} \gu{જાત}.
\end{quote}

\begin{quote}
\emph{jo have māro dāv hot to hũ jītī jāt.}
\end{quote}

\begin{quote}
If it were my turn now, I would win.
\end{quote}

A masculine noun.

\begin{cousinweb}
\textbf{\gu{દાવ}} (\emph{dāv}) is filed by CDIAL under Sanskrit \emph{dātu}, \textquotedblleft{}a share\textquotedblright{}, with the Gujarati word glossed \textquotedblleft{}a turn, a throw\textquotedblright{} (CDIAL 6258).
\end{cousinweb}

\begin{grammarlens}[title={What you've built}]
One more word for saying what happens if.
\end{grammarlens}

\subsection*{Your turn}
\begin{itemize}
\raggedright
  \item \emph{Say it:} \emph{dāv}
  \item \emph{Say it:} \emph{dāv}, once more
  \item \emph{Say it:} \emph{jo have māro dāv hot to hũ jītī jāt.} — If it were my turn now, I would win.
  \item \emph{Recall:} say the Gujarati for a die, then the Gujarati for chaupat, then say \emph{dāv} again
\end{itemize}

\begin{tcolorbox}[breakable,colback=teal!4,colframe=teal!35!black,title={Before you move on}]
What does \gu{દાવ} mean? (\textquotedblleft{}A turn (in a game)\textquotedblright{}.) Say it once more.
\end{tcolorbox}
\section[\texorpdfstring{\gu{બાજી} (bājī)}{બાજી (bājī)}]{\gu{બાજી} (bājī) — a round (of a game)}

\label{lesson:GU-C533-baji}



\begin{quote}
Before the new one: say the Gujarati for chaupat, then the Gujarati for a turn (in a game).
\end{quote}

\subsection*{You'll want to know: \gu{બાજી}}
\textbf{\gu{બાજી}} — \emph{bājī} — \textquotedblleft{}a round (of a game)\textquotedblright{}.

\begin{quote}
\gu{જો} \gu{તું} \gu{આ} \gu{બાજી} \gu{જીતીશ} \gu{તો} \gu{આપણે} \gu{બીજી} \gu{બાજી} \gu{રમીશું}.
\end{quote}

\begin{quote}
\emph{jo tũ ā bājī jītīś to āpaṇe bījī bājī ramīśũ.}
\end{quote}

\begin{quote}
If you win this round, we'll play another.
\end{quote}

A feminine noun.

\begin{cousinweb}
\textbf{\gu{બાજી}} (\emph{bājī}) is usually traced to Persian \emph{bāzī}, \textquotedblleft{}play, a game\textquotedblright{}.
\end{cousinweb}

\begin{grammarlens}[title={What you've built}]
One more word for saying what happens if.
\end{grammarlens}

\subsection*{Your turn}
\begin{itemize}
\raggedright
  \item \emph{Say it:} \emph{bājī}
  \item \emph{Say it:} \emph{bājī}, once more
  \item \emph{Say it:} \emph{jo tũ ā bājī jītīś to āpaṇe bījī bājī ramīśũ.} — If you win this round, we'll play another.
  \item \emph{Recall:} say the Gujarati for chaupat, then the Gujarati for a turn (in a game), then say \emph{bājī} again
\end{itemize}

\begin{tcolorbox}[breakable,colback=teal!4,colframe=teal!35!black,title={Before you move on}]
What does \gu{બાજી} mean? (\textquotedblleft{}A round (of a game)\textquotedblright{}.) Say it once more.
\end{tcolorbox}
\section[\texorpdfstring{\gu{ખેલવું} (khelvũ)}{ખેલવું (khelvũ)}]{\gu{ખેલવું} (khelvũ) — to play a game of chance}

\label{lesson:GU-C533-khelvu}



\begin{quote}
Before the new one: say the Gujarati for a turn (in a game), then the Gujarati for a round (of a game).
\end{quote}

\subsection*{You'll want to know: \gu{ખેલવું}}
\textbf{\gu{ખેલવું}} — \emph{khelvũ} — \textquotedblleft{}to play a game of chance\textquotedblright{}.

\begin{quote}
\gu{જો} \gu{તું} \gu{આ} \gu{દાવ} \gu{ખેલત} \gu{તો} \gu{તું} \gu{જીતી} \gu{જાત}.
\end{quote}

\begin{quote}
\emph{jo tũ ā dāv khelat to tũ jītī jāt.}
\end{quote}

\begin{quote}
If you had played that turn, you would have won.
\end{quote}

Past, with \gu{તેણે} (he or she …): \gu{તેણે} \gu{દાવ} \gu{ખેલ્યો}, \emph{teṇe dāv khelyo}.

\begin{cousinweb}
\textbf{\gu{ખેલવું}} (\emph{khelvũ}) comes from \emph{*khēḍ}, \textquotedblleft{}play\textquotedblright{}, a form CDIAL reconstructs, and CDIAL lists the Gujarati verb (CDIAL 3918). Beside \textbf{\gu{રમવું}} it is the bolder word, for a game played for a stake.
\end{cousinweb}

\begin{grammarlens}[title={What you've built}]
One more word for saying what happens if.
\end{grammarlens}

\subsection*{Your turn}
\begin{itemize}
\raggedright
  \item \emph{Say it:} \emph{khelvũ}
  \item \emph{Say it:} \emph{khelvũ}, once more
  \item \emph{Say it:} \emph{jo tũ ā dāv khelat to tũ jītī jāt.} — If you had played that turn, you would have won.
  \item \emph{Recall:} say the Gujarati for a turn (in a game), then the Gujarati for a round (of a game), then say \emph{khelvũ} again
\end{itemize}

\begin{tcolorbox}[breakable,colback=teal!4,colframe=teal!35!black,title={Before you move on}]
What does \gu{ખેલવું} mean? (\textquotedblleft{}To play a game of chance\textquotedblright{}.) Say it once more.
\end{tcolorbox}
